% !TeX root = ../papers/nc-p3.tex
\section{正次元の分岐中心は直線に限る}
\label{sec:branching-lines-only}
\begin{theorem}\label{thm:branching-lines-only}
標準正則四窓の各射影で、直線ファイバーを持つ点の被約軌跡は有限個の点と有限本の直線の和である。
各直線が次の文字に掃く曲面は平面、滑らかな二次曲面、または二次錐である。
滑らかな二次曲面を掃く場合には、その直線上の全点は同じ一つの前点を持つ。
またいずれかの辺の射影像が射影三空間全体なら、全窓は添字付き可換多項式環の一類で尽きる。
\end{theorem}
\begin{proof}[三次元像と曲面次数の上界]
一般点で関係の評価階数が3なら、その辺から後続点への有理写像がある。
一般ファイバーが正次元なら逆方向の直線ファイバーの中心は曲面以上を占める
（全ファイバーは高々直線）。前節に反するので像は三次元、逆も一般一点で双有理である。
座標を共通因子のない四つの独立な同次数式 $f_i$ で表す。
六関係より $f_i$ に全四一次式を掛けた包の次元は10以下である。
このとき $f_i$ は一次式である。実際、基底を消去して異なる四つの先頭単項式を取る。
その各四倍のうち、二単項式間の一致は高々一つである。従って相異なる積は $16-6=10$ 以上で、
等号の場合全対に一致があり三重の一致はない。四単項式は同じ単項式 $h$ に四文字を掛けたものである。
最初の二つを $h x_i,h x_j$ とすると、両方と一文字倍で会う第三のものは
$h$ の倍か $h x_i x_j$ の約数である（指数の差が一文字の交換であることによる）。
後者は三重一致になるからである。各対のS多項式は一つ上の次数で零に簡約できる。
この次数の先頭単項式は既に十個の倍で尽きるためである。
従って四式はGröbner基底であり、$h$ が非定数なら初期イデアルの高さは1である。
元のイデアルも高さ1なら四式に共通因子があり矛盾する。この証明は四次検査による延長を使わない。
従って双有理写像は射影線型で、そのグラフを消す六関係は交代関係全体に移される。
一般の零点対では隣の評価も階数3以上であり、直交不等式から3になるので、隣接辺にも同じ議論を適用できる。
全整数位置で順に基底を運ぶと全関係が可換関係となる。

以後この場合を除く。各辺で一般階数は4である。
四次小行列式の最大公約因子は次数2以下である。仮に次数4なら四行六列の評価行列の核二平面は一般に一定で、
非零関係が全ての点で消えてしまう。次数3なら、小行列式を割ったPlücker写像の像の閉包は線型空間である。
全核に共通の関係はないため定点や直線ではなく、前節と同じ交わる直線族の分類から
核は全て $R$ の一定の三平面 $T$ に入る。すると $T$ の評価は全点で階数1以下である。
一次行列の階数1以下の線型空間は共通像直線または共通余像直線を持つので、
三つの独立な関係が同じ末尾文字を持つか、非零関係が恒等零となる。どちらも不可能である。
従って各射影で全曲面成分の次数の和は2以下である。
\end{proof}

\begin{proof}[曲線と二項表示]
分岐中心の既約曲線を $C$ とする。その後続直線は相異なり既約曲面を掃き、
一般の曲面点には前点が一つしかない。曲面が平面なら、双対平面で後続直線を表した曲線の次数は1である。
実際、その次数が一般の平面点を通る直線の本数である。よって直線束の中心点を共有する。
二次錐の母線も頂点を共有する。共通点があれば、その前点は一つの直線に含まれるため $C$ 自身が直線である。
残りでは後続は滑らかな二次曲面の一方の母線族全部である。$C$ からそのパラメータ直線への射は
単射で、標数零なので有限双有理、標的が正規なので同型である。
すなわち $C=\mathbb P^1$ と書け、後続点は二項座標で $s\otimes t$、$s\in C$、$t$ は自由である。

非線型な $C$ を仮定する。平面に入れば滑らかな有理平面曲線なので二次曲線である。
非平面の場合に $C$ の文字の制限を $A\subset S^d$、$S^j=H^0(\mathbb P^1,\mathcal O(j))$ とする。
$A$ は基点自由な四空間である。この辺の六関係は、$t$ の二座標ごとに
$A\otimes S^1\to S^{d+1}$ の乗法核に入る。従ってこの核は少なくとも三次元であり、
$\dim S^1 A\le5$ である。任意の非零二項形式空間 $B$ で $\dim S^1B\ge\dim B+1$ である
（等次元なら二一次式の非定数比が有限次元空間を保ってしまう）。
等号なら $B=h S^{\dim B-1}$ である。実際、$x b=y b'$ の解を
$b=y c,b'=x c$ と書くと、その $c$ の空間は $\dim B-1$ 次元で、その一次式倍は全て $B$ に入り
同じ最小増加を持つ。次元で帰納すればよい。よって $A=S^3$ で $C$ は有理正規三次である。

この二つを $d=2,3$ とし、射影平面または三空間への文字制限は完全系 $S^d$ である。
位置を $0=C,1=(s,t),2,3$ とする。位置0を $C$ の線型包へ制限した関係 $R'_0$ は
\[
        [x,y]S^{d-1}(x)\otimes S^1(t)
\]
に等しい。$x,y$ は位置0と位置1の母線パラメータの二項変数であり、括弧は二行列式である。
$d=3$ では次元6で等号、$d=2$ ではこの包の次元は4で、平面を消す一文字 $a$ に始まる関係は高々二つなので
やはり等号である。後者では $a$ に始まる純関係があり $a$ に終わる純関係は禁じられる。
従って $R_3$ の位置0への同じ射影 $R'_3$ は六次元のままである。$d=3$ でも当然六次元である。
\end{proof}

\begin{proof}[前後に直線を持つ曲線]
$C$ の一般点の逆方向の評価階数は2か3である。2の場合は $C$ 全体で2となり、
逆順にも同じ曲面の議論が適用できる。$C$ を非線型としているので、前点側も滑らかな二次曲面の
一方の母線族である。その母線パラメータも $C$ からの射影同型で同定し、変数を $v$、
自由なもう一方の二項変数を $u$ とする。
制限された四テンソルは $R'_0,R'_3$ に入るので
\[
              [x,y][x,v]h(x,t,z,u),\qquad \deg_x h=d-2
\]
と書ける（$z$ は位置2の四文字）。二つの括弧は互いに素である。
$d=2$ なら $y,t,z\mid v,u,x$ の階数は高々 $2\cdot2=4$ である。
実際、二つの括弧のこの分割での階数は高々2、残る $h$ の $u$ 側も二次元である。
しかし同じ切断の像は六次元である。位置0だけを制限しており、$R_3$ の六次元像が保たれるからである。
$d=3$ では $s\in C$ を固定すると、その両端支持で文字を割ることは $y=v=s$ での評価である。
系\ref{cor:endpoint-line-line-cut}により残る $[x,s]^2h(x,t,z,u)$ の位置0対他の階数は高々1である。
従って $h$ は $x$ の固定一次式の倍であり、$S^3(x)$ の全支持に反する。
\end{proof}

\begin{proof}[一般前点が唯一の場合]
前点の評価階数が一般に3とする。制限四テンソルを
$[x,y]f(x,t,z,u)$ と書き、今の $u$ は位置3の四文字とする。
$u$ に値を持つ次数 $d-1$ の二項形式の空間 $A'$ を、$t,z$ による $f$ の全縮約の包とする。
$R'_3=S^1 A'$ なので $\dim S^1 A'=6$ である。一般の $s$ での $A'(s)$ の次元は
端支持の階数3である。$s$ の消滅一次式を掛けた $A'$ は評価核に入るので
$6\ge\dim A'+3$ から $\dim A'=3$ である。そこで
\[
               f=\sum_{j=1}^3 K_j(t,z) A_j(x,u)
\]
と分解すると、$K_j$ は独立で、$R_1=S^1(y)\otimes\langle K_j\rangle$ である。
実際、制限後の他側切断の像はこの六次元空間に入り、$R'_3$ の階数により等号である。
制限の前の $R_1$ も六次元だから変わらない。

位置0を $S^d$ に割った商で、位置1の二項二文字 $y,t$ を一旦二つの連続する位置と読む。
$R'_0$ は $x,y$ だけへの関係と自由な $t$ の積、$R_1$ は自由な $y$ と $t,z$ だけへの関係の積である。
従って $y$ と次の $t$ の間には関係がなく、一期を $t,z,u,x,y$ としたブロックが独立に繰り返される。
厳密には元の経路を $z$ から始め $x$ で終え、内部に必要回数の完全ブロックを入れればよい。
端の部分ブロックは完全ブロックに非零商があれば非零である。$d=3$ ではこの分離は元の代数そのもので、
完全ブロックの次元が0,1なら有界、2以上なら指数成長となり標準Hilbert関数に反する。

$d=2$ の商では次元が2以上であることを示す。まず $t,z$ の二文字の商は $8-3=5$ 次元である。
ブロックの右端 $x,y$ の商は乗法で $S^3$ になり、$u,x,y$ の商は
$(V_3\otimes S^3)/(S^2 A')$ なので $16-9=7$ 次元である。
実際、三つの $A_j$ は一般点で独立なため $S^2$ 係数の間にも関係はない。
$R_2$ の六つのうち四つは $A_j=x_0 A_{0j}+x_1 A_{1j}$、
$K_j=t_0 K_{0j}+t_1 K_{1j}$ として
$H_{bi}=\sum_j K_{bj}(z) A_{ij}(u)$ ($b,i=0,1$) である。
これらは制限四テンソルの位置0,1からの縮約で、$R'_0$ の階数4から独立である。

$5\cdot7$ 次元の積へ残る中央関係を挿入する。
$H_{bi}$ の全挿入の像は高々 $3\cdot5=15$ 次元である。
左側で $t_c K_{bj}$ の四つの添字組には全 $j$ で
$\sum_b t_b K_{bj}=0$ がある。右側で $A_{ij}\otimes p_i$、$p_i\in S^3$ の二組には
$(p_0,p_1)=(x_0g,x_1g)$, $g\in S^2$ という三次元の核が全 $j$ に共通である。
従ってこの挿入写像は左右で次元 $4-1$ と $8-3$ の商の積を経由する。
追加二関係の挿入は $2\cdot2\cdot4=16$ 次元以下で、完全ブロックは $35-15-16=4$ 次元以上である。
この商の指数増大も元の標準成長に反する。これで全非線型曲線を除いた。

最後に $C$ が直線で後続が滑らかな二次曲面の母線のとき、位置0の二文字への射影では
$R'_0=[x,y]S^1(t)$ である。包含はグラフでの消滅、等号は各点の端支持2による。
四テンソルは $[x,y]f(t,z,u)$ に制限されるので、位置0の任意の点で逆端の支持は
$f$ の $u$ 支持に等しく一定である。その次元は2以上、一般点で3以下であり、
2なら異なる点が同じ逆支持平面を持つ。従って3で、同じ唯一の前点が全直線に対して決まる。
\end{proof}
